\SourceLeaf{88}%

\noindent sodann

\[ T_2 = \frac{1}{2}\mathord{.}\frac{n}{n+\lambda}\frac{d\Theta_1}{dt} = \frac{n(n+2\lambda-1)}{2}\, t^{n+2\lambda-2} \]

\noindent und daraus auf ähnliche Weise

\[ \Theta_2 = \frac{(n+\lambda)(n+3\lambda-1)}{2}\, t^{n+3\lambda-2}, \]
\[ T_3 = \frac{1}{3}\frac{n}{n+\lambda}\frac{d\Theta_2}{dt} = \frac{n(n+3\lambda-1)(n+3\lambda-2)}{2.3}\, t^{n+3\lambda-3}, \]
\[ \Theta_3 = \frac{(n+\lambda)(n+4\lambda-1)(n+4\lambda-2)}{2.3}\, t^{n+4\lambda-3}, \]
\[ T_4 = \frac{1}{4}\frac{n}{n+\lambda}\frac{d\Theta_3}{dt} = \frac{n(n+4\lambda-1)(n+4\lambda-2)(n+4\lambda-3)}{2.3.4}\, t^{n+4\lambda-4}, \]

\noindent wobei das Bildungsgesetz klar ist.

\Needspace{8\baselineskip}
\begin{center}\textls[90]{II. Allgemein.}\end{center}

\noindent Grundgleichung: \quad $x = t + u\,\psi x,$

\noindent Gesucht: \quad $fx = ft + T_1 u + T_2 uu + T_3 u^3 + \text{u.\,s.\,w.};$

\noindent $\psi x$ und $fx$ sind gegebene Funktionen von $x$, während $T_1$, $T_2$, $T_3$ usw. noch zu bestimmende Funktionen von $t$ sind. Ihre Bestimmung hängt von der Beschaffenheit der Funktionen $\psi x, fx$ ab. Die Ableitung $\dfrac{dfx}{dx}$ soll mit $Fx$ bezeichnet werden.

Sei nun $\varphi x$ eine andere Funktion von $x$ und $\Phi x = \dfrac{d\varphi x}{dx}$ ihre Ableitung. Somit kann man auch schreiben:

\[ \varphi x = \varphi t + \Theta_1 u + \Theta_2 uu + \Theta_3 u^3 + \text{u.\,s.\,w.}, \]

\noindent wobei $\Theta_1, \Theta_2, \Theta_3$ usw. ebenfalls Funktionen von $t$ sind, die von der Natur der Funktionen $\psi, \varphi$ ebenso abhängen wie $T_1, T_2, T_3$ usw. von $\psi, f$. Sollte sich zeigen, dass in der letzteren Abhängigkeit die Funktion $f$ selbst nicht vorkommt, sondern nur $F$, so gehen, da $\psi$ in beiden dieselbe Rolle spielt, $T_1, T_2, T_3$ usw. in $\Theta_1, \Theta_2, \Theta_3$ usw. über, wenn man im Ausdruck für die erstere Funktion überall dort, wo $Ft$ vorkommt, dafür $\Phi t$ einsetzt.
